\problem{Gaz parfait ultrarelativiste dans le régime classique}

\question{
    Calculez la fonction de partition d'une particule ultrarelativiste se déplaçant en trois dimension (3D) et en absence de champ magnétique
    }

Pour ce problème, on approxime $\epsilon \approx c\abs{\vb{p}}$.
On peut trouver la fonction de partition d'une particule ultralativiste tel que:
\[
    \zeta &= \sum_{\va{k},~\va{p},~s} e^{-\beta \epsilon_{\va{k},~\va{p},~s}} \\
    &= \int_{0}^{\infty} \frac{\dd[3]r \dd[3]p}{h^3} e^{-\beta \epsilon_{\va{k},~\va{p}}} \sum_s e^{-\beta \epsilon_s}. \notag 
    \intertext{Posons $\sum_s e^{-\beta \epsilon_s} = g$, le facteur de dégénéressence du spin. On obtient alors}
    \zeta &= \frac{gV}{h^3} \int_{0}^{\infty} \dd[3]p e^{-\beta c \abs{\va{p}}}. \notag \\
    \intertext{En résolvant l'intégrale on trouve} 
    \zeta &= \frac{gV}{h^3} 8\pi \notag \\
    &= \frac{gV}{\lambda_{th}^3},~~~ \text{où $\lambda_{th} = \frac{h\beta c}{\qty(8\pi)^{\frac{1}{3}}}$.}
\]

\question{
    Montrez que votre définition précédente de $\lambda_{th}$ pour les particules ultrarelativistes est cohérente (modulo un facteur numérique de l'ordre d'unité) avec la valeur moyenne de la longueur d'onde de de Broglie.
    }

On peut approximer la moyenne de la longueur d'onde de de Broglie tel que
\[
    \expval{\lambda_{th}} &= \expval{\frac{h}{p}} \notag \\
    &\approx \frac{h}{\expval{\abs{\va{p}}}}.
\]

Trouvons $\expval{p}$:
\[
    \expval{p} &= \frac{\int_{0}^{\infty} \dd[3] p \abs{\va{p}} e^{-\beta c \abs{\va{p}}}}{\int_{0}^{\infty} \dd[3] p' e^{-\beta c \abs{\va{p'}}}}. \\
    &= \frac{6}{\qty(\beta c)^4} \frac{\qty(\beta c)^3}{2} \notag \\
    &= \frac{3}{\beta c}.
\]

On obtient alors
\[
    \expval{\lambda_{th}} &\approx \frac{hc}{3k_B T}
\]
ce qui est cohérent avec $\lambda_{th}$ puisque
\[
    \lambda{th} &= \qty(8\pi)^{\frac{1}{3}} \notag \\
    &= \frac{3}{\qty(8\pi)^{\frac{1}{3}}} \expval{\lambda_{th}} \notag \\
    \Rightarrow \expval{\lambda_{th}} &\propto \lambda_{th}.
\]

\question{
    Calculez le potentiel chimique, l'énergie moyenne, la pression et l'entropie pour un gaz classique ultrarelativiste en 3D et en absence de champ magnétique. Indiquez les similarités et les différences par rapport aux propriétés du gaz classique non relativiste étudiées en classe.
    }

Pour un gaz classique, la fonction de partition est définie tel que:
\[
    Z_N &= \frac{\zeta^N}{N!} \notag \\
    \Rightarrow \ln\qty[Z_n] &\approx N\ln\qty[\frac{gV}{N\lambda_{th}^3}]+N.
\]

Le potentiel chmique est défini tel que:
\[
    \mu &= \eval{\pdv{F}{N}}_{T,~V} \\
    \intertext{avec $F = -k_BT\ln\qty[Z_n]$.}
    \Rightarrow \mu &= -k_BT \pdv{N}\eval(N\ln\qty[\frac{gV}{N\lambda_{th}^3}] + N|_{T,~V} \notag \\
    &= k_B t \ln\qty[\frac{n \lambda_{th}^3}{g}],~~~ \text{où $n = \frac{N}{V}$.}
\]

L'énergie moyenne est définie telle que:
\[
    \expval{\epsilon} &= -\pdv{\ln Z_N}{\beta} \notag \\
    &= 3NkBT.
\]

La pression est définie telle que:
\[
    p &= k_BT \eval{\pdv{\ln Z_N}{V}}_{T,~N} \notag \\
    &= Nk_BT.
\]

L'entropie est définie telle que:
\[
    S &= -\eval{\pdv{F}{T}}_{V,~N}. \notag
\]
Or, l'énergie libre de Helmholtz est
\[
    F &= \expval{\epsilon} - TS \notag \\
    \iff S &= \frac{\expval{\epsilon} - F}{T}. \notag \\
    \intertext{Or, on connait $\expval{\epsilon}$ et $F$, d'où}
    S &= 3Nk_B \qty(N\ln\qty[\frac{gV}{N\lambda_{th}^3}] + 4).
\]

Le gaz relativiste et le gaz non-relativiste ont comme similarité leur la forme du potentiel chimique ($\mu \propto k_B T \ln\qty[\frac{n\lambda_{th}^3}{g}]$) et la structure de leur fonction de partition ($\zeta \propto \frac{gV}{\lambda_{th}^3}$).

Ils diffèrent de par leur longueur d'onde thermique ($\lambda_{th}^{NR} \propto T^{-\frac{1}{2}}$ contre $\lambda_{th}^{UR} \propto T^{-1}$), leur énergie moyenne ($\expval{\epsilon_{NR}} = \frac{3}{2}Nk_BT$ contre $\expval{\epsilon_{UR}} = 3Nk_BT$) et de leur relation pression-énergie ($p_{NR} = \frac{2\expval{\epsilon_{NR}}}{3V}$ contre $p_{UR} = \frac{\expval{\epsilon_{UR}}}{3V}$).

\clearpage

